\chapter{Virando em Direcção ao Vazio}

\lettrine{A}{través da reflexão, trazemos à consciência} o estado das condições como elas são
agora. Tendo nascido, estamos agora nesta idade, sentindo"-nos assim, neste
momento e neste lugar. É a realidade como é. Isso não pode ser mudado por nós. É
simplesmente a inevitabilidade do nascimento que agora é a realidade como é.

E com essa reflexão, obtemos uma perspectiva da realidade como é, em vez da
reacção à realidade como é. Se não reflectirmos, simplesmente reagimos à
realidade como é.

Se nos sentimos felizes, ficamos eufóricos ``quero ser um monge para o resto da
minha vida e dedicar"-me ao Dhamma. Dhamma é o caminho para mim. O único caminho,
o caminho verdadeiro\ldots{}'' E saímos e aborrecemos as pessoas com uma arenga
acerca da importância do budismo no mundo, porque estamos eufóricos e
sentimo"-nos positivos e confiantes. Até esse sentimento de estarmos inspirados e
confiantes, cheios de fé, devoção e todo esse tipo de coisas -- é a realidade
como é. Podemos sentir muita fé e confiança no que estamos a fazer.

Ou podemos sentir o oposto: perde"-se a fé e sentimos que é uma perda de tempo
``Desperdicei a minha vida. Não tem valor, não cheguei a lado nenhum. Não surtiu
nada para mim. Já não acredito mais nisto, estou farto disto.'' Ou podemos
sentir indiferença: ``Está tudo bem, não sei o que fazer mais. Melhor do que
trabalhar numa fábrica.'' Se é assim que nos sentimos neste momento -- seja no
extremo ou simplesmente indiferença -- é a realidade como é.

Portanto, apenas observemos quando estamos a sentir uma tremenda energia e a
sentir"-nos positivos, ou quando há falta dela e estamos muito críticos. Quando
estamos deprimidos e quando não nos estamos a sentir muito bem, ou quando
estamos cansados, é difícil elevar o sentimento inspirado. Nessas
circunstâncias, temos tendência para muito rapidamente reparar no que há de
errado com as coisas. A forma como alguém atravessa uma sala pode irritar"-nos.
Alguém assoa o nariz com muita força -- e oh, isso é nojento! Mas quando nos
estamos a sentir cheios de inspiração e devoção, não nos importamos simplesmente
com as falhas disto ou daquilo, apenas nos prendemos por este sentimento de
devoção e fé. Estas percepções são para serem reflectidas como a realidade que é
agora. Tem que ser assim, porque não pode ser de mais forma alguma neste
momento. Nós sentimo"-nos assim, sentimo"-nos cansados ou revigorados, ou o que
for -- é a realidade como é.

Estes são os resultados de termos nascido e de viver as nossas vidas e de
estarmos sujeitos às mudanças nas condições da sensualidade. Então
apercebamo"-nos. Apercebamo"-nos mesmo, o que adicionamos às condições existentes.
Nas meditações que duram a noite inteira, podemos sentir sono ou cansaço;
apercebamo"-nos do que colocamos em cima desse sentimento. Observemos a sensação
em si, mas mantenhamos uma postura, em vez de apenas reagirmos à sensação de
cansaço, tentando aniquilar a sensação ao segui"-la e mergulhando na letargia.

Quando estivermos realmente convencidos de que estamos tão cansados que não há
nada que possamos fazer a esse respeito, que mesmo endireitando o nosso corpo
parece algo totalmente impossível, seguramo"-lo direito por um período de tempo.
Apercebemo"-nos e aprendemos quanta energia é necessária para manter um corpo
em~pé.

Quanta energia é necessária para interromper o processo do pensamento? Já alguma
vez repararam nisso? ``Simplesmente não consigo parar de pensar'' continua a
mente indefinidamente. ``Não consigo parar, o que posso fazer?'', ``Não sei como
parar de pensar'' continua. ``Eu não consigo parar\ldots{}'' Eu conheço a
condição porque sempre tive este problema com a minha mente que parecia estar
interminavelmente a pensar em alguma coisa. E o desejo de parar de pensar e o
esforço para nos vermos livres disso, cria as condições para mais pensamento!

É preciso esforço para o fazer, não apenas pensar nisso. Lembro"-me de uma vez,
um australiano que era fanático sobre o Abhidhamma veio a Wat Pah Pong. Este
homem tinha uma missão -- quando os ocidentais se metem no Abhidhamma, tornam"-se
como cristãos"-nascidos"-de-novo -- mas ele não sabia meditar; não acreditava que
a meditação funcionasse e descobriu tudo com os seus conceitos de Abhidhamma.
Ele pensava que não conseguimos parar de pensar. Então ele disse, ``Você está
sempre a pensar e não consegue parar de pensar''. E eu disse, ``Mas você pode
parar de pensar''. E ele disse, ``Não, não posso\ldots{}'', e eu disse,
``acabei de parar de pensar\ldots{}'', e ele disse, ``Não, não parou''!

É inútil continuar a falar com alguém assim. É necessário estarmos alerta para
sabermos quando não estamos a pensar, então usamos um pensamento como ``Eu não
consigo parar de pensar'' e deliberadamente pensamos isso. Isto foi o que eu
fiz, porque eu era habitualmente um pensador obsessivo. Portanto, em vez de
tentarmos parar de pensar, se formos adversos a isso, vamos então para o extremo
oposto e pensamos deliberadamente algo.

E observemo"-nos a pensar deliberadamente, de forma que não seja só um vago
processo de pensamento em que a nossa mente não para de girar às voltas.

Usemos a faculdade da nossa sabedoria; deliberadamente pensemos algo, algum
pensamento que seja totalmente neutro e desinteressante como ``eu sou um ser
humano''. Então pensamos, mas observamos o espaço antes de pensarmos e então
dizemos deliberadamente ``Eu sou um ser humano''. Logo notamos o fim do que
acabámos de dizer, no momento em que paramos de pensar. Se prestarmos atenção no
antes e no depois do pensamento, em vez de no pensamento em si, mantenhamos essa
atenção onde não há pensamento. Investiguemos o espaço ao redor do pensamento, o
espaço onde o pensamento vem e vai, ao invés de pensarmos.

Então estamos cientes de uma mente vazia onde há apenas consciência, mas nenhum
pensamento. Isto pode durar apenas um segundo porque começamos a querer pensar.
Então só precisamos de continuar a ser mais conscientes pensando novamente em
algo. Com a prática, podemos até usar pensamentos muito desagradáveis. Por
exemplo, podemos ter fortes sentimentos emocionais como ``eu não presto, sou
inútil'' e isso pode ser uma obsessão. Na mente de algumas pessoas, pode
tornar"-se num pano de fundo em suas vidas. Então podem tentar pensar: ``Eu não
devia pensar assim. O Venerável Sumedho diz que sou bom. Mas sei que não
presto.'' No entanto, se seguirmos essa obsessão e usarmo"-la como um pensamento
consciente: ``Eu não presto'', começaremos a ver o espaço em seu redor. E já não
soa tão absoluto não é verdade? Quando se torna obsessivo, soa absoluto,
infalível, a verdade honesta, a verdade real: ``Isto é o que eu realmente sou,
não presto.'' Mas quando realmente retiramos isso do contexto da obsessão -- e
pensamos nisso deliberadamente -- observamos isso objectivamente.

Esse sentimento de `eu' e `meu' é apenas um hábito da mente; não é a verdade. Se
realmente reflectirmos no `eu/eu sou' e olharmos para isso objectivamente, esse
sentimento criado por esse `eu sou' e `eu sou assim' ou `eu deveria ser/não
deveria ser' é muito diferente de quando estamos apenas a reagir.

Ao contemplar as Quatro Nobres Verdades, temos a verdade do sofrimento; o seu
surgimento; a sua cessação e então o Caminho. Não é possível conhecer o Caminho
e a saída do sofrimento até que estejamos cientes de onde tudo cessa -- na
própria mente. A mente ainda está bem viva e alerta mesmo quando não há
pensamento nela; mas se não percebermos isto, então acreditamos que estamos
sempre a pensar. É assim que parece.

Nós só concebemos uma ideia de nós mesmos quando estamos a pensar, porque
estamos identificados com a memória e sentimento de ``eu sou'' ou ``eu não
sou''. Esse ``eu mesmo'' é muito uma percepção condicionada e programada na
mente. Enquanto acreditarmos nesta percepção e nunca a questionarmos, vamos
sempre acreditar que somos uns pensadores obsessivos\ldots{} e que não devíamos
ser assim ou não devíamos sentir assim e que não devíamos nos preocupar -- mas
preocupamo"-nos, e somos um caso perdido e por aí adiante de um juízo a outro.

Assim, o `eu sou' é apenas uma percepção realmente -- surge e cessa na mente.
Quando cessar, observemos a cessação do pensamento. Tornemos essa cessação, essa
mente vazia, num ``sinal'' em vez de apenas criarmos mais coisas no vazio.
Podemos alcançar estados subtis de consciência fixando"-nos em objectos subtis --
como nas práticas de meditação \emph{samatha} onde se enfatiza o acalmar da
mente -- mas com a contemplação das Nobres Verdades, estamos a usar a faculdade
da sabedoria para perceber onde tudo cessa.

E, no entanto, quando a mente está vazia, os sentidos ainda estão a funcionar
bem, de forma normal. Não é como estar em transe, totalmente alheios a tudo; a
nossa mente está aberta, vazia -- ou podemos dizer inteira, completa, clara.
Então podemos trazer qualquer coisa\ldots{} como um pensamento cheio de medo.
Seguramos e pensamos deliberadamente nisso e olhamos para isso como apenas mais
uma condição da mente, ao invés de um problema psicológico. Surge e cessa; nada
existe nisso, nada em que pensamento for. É apenas um movimento na mente e por
conseguinte não é uma pessoa. Tornamo"-lo pessoal ao nos apegarmos a esse
conceito, acreditando: ``E eu sou um caso tão desesperado, sei que nunca me
poderei iluminar depois de todas as coisas que fiz; as coisas estúpidas. E eu
sou tão egoísta e cometi tantos erros. Sei que não há esperança para mim.'' Tudo
isso surge e cessa na mente!

Acreditar é ambicionar saber, não é? ``Sei o que sou e sei que não presto.''
Acreditamos nisso, e isso é precisamente a ambição do saber. Criamos essa crença
e então a mente avança nessa direcção, logo podendo encontrar todos os tipos de
provas de que nós não prestamos -- podemos até começar a ficar paranóicos:
``todas as pessoas sabem igualmente que também não presto. O Venerável Sucitto
ontem, passou por aqui e eu simplesmente soube que ele sabe que eu não presto.
Então, hoje de manhã, cheguei à sala e a irmã Rochana olhou"-me de uma forma um
pouco estranha -- ela sabe!''

Então, pela crença, podemos ver e interpretar tudo o que as pessoas fazem de uma
forma pessoal, como se todos nos estivessem a julgar e a condenar -- isso é
paranóia, não é?

Mesmo os mais belos pensamentos e aspirações, assim como os mais desagradáveis e
maléficos, surgem e cessam na mente. Agora, não me compreendam mal -- não estou
a dizer que pensamentos bons e maus são a mesma coisa. Eles têm a mesma
característica de surgir e cessar, é só isso. Noutros aspectos são diferentes.
Bons pensamentos são bons pensamentos, maus pensamentos são maus pensamentos!
Portanto, não estou a dizer que não há problema em ter maus pensamentos, mas sim
a apontar para além da qualidade do pensamento: amor e ódio surgem e cessam na
mente. Nesta perspectiva, dirigimo"-nos para a mente reflexiva para a qual a
maioria das pessoas está totalmente inconsciente. As pessoas geralmente só estão
conscientes de si mesmas como personalidade, uma emoção ou um pensamento -- em
outras palavras, como uma condição.

Para praticar, não nos preocupemos com as qualidades que passam pela mente: quão
maravilhosas, interessantes, bonitas, feias, desagradáveis ou neutras elas podem
ser. Não estamos a investigar qualidades ou a negar a qualidade de qualquer
pensamento, mas somente a apercebermo"-nos da realidade como é. Então
simplesmente deixamos o pensamento sozinho de maneira a que ele cesse. Criamos
um pensamento, colocamo"-lo na mente deliberadamente e deixamo"-lo ir. Deixar ir
não significa afastá"-lo: deixamos sim o pensamento sozinho, estamos conscientes
disso o tempo todo; do momento anterior ao pensamento, dos interstícios e do
final.

O espaço ao redor do pensamento -- nós não nos percebemos muito disso, certo? É
como o espaço nesta sala, em que tenho que enfatizar a atenção nisso para nos
apercebermos dele. Ora, o que é que é preciso para estarmos cientes do espaço
nesta sala? Temos que estar alerta. Com os objectos da sala não precisamos estar
alerta, podemos apenas sentir"-nos atraídos ou repelidos: ``Eu não gosto disso,
eu gosto daquilo.'' Podemos simplesmente reagir à qualidade da beleza e do que é
feio, quer isso nos agrade ou desagrade. É o nosso hábito, certo? Na nossa vida
inclinamo"-nos normalmente a reagir ao prazer e à dor, à beleza e ao que é feio.
Então olhamos para a beleza e dizemos: ``Oh, olhe para isto! Não é absolutamente
fantástico'' ou pensamos ``Oh, que nojento!''

Mas os objectos bonitos e feios estão todos no espaço e para se perceber o
espaço, temos que retirar a atenção tanto do que é belo como do que é feio.
Claro que os objectos ainda lá estão, não precisamos deitá"-los fora; não
precisamos deitar o prédio abaixo para que possamos ter um espaço aqui. Mas se
não concentrarmos amor ou ódio no que está naquela sala, se não fizermos nenhum
juízo apreciativo, a nossa atenção retira"-se dos objectos e apercebemo"-nos do
espaço.

Então, temos uma perspectiva do espaço numa sala como esta. Pode reflectir sobre
isso. Qualquer pessoa pode entrar e sair neste espaço. A mais bonita, a mais
feia, o santo e o pecador, podem ir e vir neste espaço e o espaço nunca é
prejudicado, arruinado ou destruído pelos objectos que vão e vêm neste espaço.

A mente funciona no mesmo princípio. Mas se não temos o hábito de observar o
espaço da nossa mente, não estamos cientes do espaço que a mente realmente é.
Então não temos consciência do vazio da mente, porque estamos sempre apegados a
uma ideia, opinião ou temperamento.

Na meditação com compreensão lúcida, acabamos por reflectir sobre os cinco
\emph{khandhas} -- no corpo ``\emph{rūpa}'', sentimentos ``\emph{vedanā}'',
percepção ``\emph{saññā}'', formações mentais ``\emph{saṅkhārā}'' e cognição
``\emph{viññāṇa}''. Podemos querer nos livrar deles, mas isso é outra condição,
outro \emph{saṅkhārā} que criamos. Então nós investigamo"-los até que não nos
iludam mais e permitimos que eles cessem na mente vazia. Quando pensamos ``o meu
corpo ainda está aqui -- como ele cessa? Ainda está aqui, não está?''
Consideremos que o corpo viverá a duração da sua vida desde que nasceu e
desaparecerá quando sua força cármica terminar.

O que aconteceu com Napoleão? O que aconteceu com a Rainha de Sabá? E Confúcio e
Lao Tzu e Maria Antonieta, Beethoven e Bach? São memórias nas nossas mentes,
apenas percepções no agora na mente das pessoas. Mas na realidade era isso que
elas sempre foram, mesmo quando os seus corpos estavam vivos!

``Venerável Sumedho'' é uma percepção na mente -- é uma percepção na minha
mente, é uma percepção na vossa mente. No momento, a percepção disso é: ``o
Venerável Sumedho está vivo e activo.'' Quando o corpo morre, então a percepção
muda para ``o Venerável Sumedho está morto''. É só isso, não é? A percepção da
morte está lá a acompanhar o nome Sumedho, que agora está vivo e activo. Assim,
à medida que vivenciamos o corpo, esta é uma percepção que surge na mente -- e
que cessa na mente vazia.

Com esta realização da mente vazia, podemos desenvolver o Nobre Óctuplo Caminho
de uma forma muito sábia. O Nobre Óctuplo Caminho baseia"-se na Compreensão
Correcta, e isso é a compreensão da cessação.

\photoPage{image-010-two-siladhara.jpg}{Sīladharas em Hammer Wood, Chithurst}{Foto por Mike Holmes}
